% Options for packages loaded elsewhere
\PassOptionsToPackage{unicode}{hyperref}
\PassOptionsToPackage{hyphens}{url}
\PassOptionsToPackage{dvipsnames,svgnames,x11names}{xcolor}
%
\documentclass[
  11pt,
  a4paper,
]{ltjsarticle}
\usepackage{amsmath,amssymb}
\usepackage{iftex}
\ifPDFTeX
  \usepackage[T1]{fontenc}
  \usepackage[utf8]{inputenc}
  \usepackage{textcomp} % provide euro and other symbols
\else % if luatex or xetex
  \usepackage{unicode-math} % this also loads fontspec
  \defaultfontfeatures{Scale=MatchLowercase}
  \defaultfontfeatures[\rmfamily]{Ligatures=TeX,Scale=1}
\fi
\usepackage{lmodern}
\ifPDFTeX\else
  % xetex/luatex font selection
\fi
% Use upquote if available, for straight quotes in verbatim environments
\IfFileExists{upquote.sty}{\usepackage{upquote}}{}
\IfFileExists{microtype.sty}{% use microtype if available
  \usepackage[]{microtype}
  \UseMicrotypeSet[protrusion]{basicmath} % disable protrusion for tt fonts
}{}
\makeatletter
\@ifundefined{KOMAClassName}{% if non-KOMA class
  \IfFileExists{parskip.sty}{%
    \usepackage{parskip}
  }{% else
    \setlength{\parindent}{0pt}
    \setlength{\parskip}{6pt plus 2pt minus 1pt}}
}{% if KOMA class
  \KOMAoptions{parskip=half}}
\makeatother
\usepackage{xcolor}
\usepackage[margin=24mm]{geometry}
\usepackage{color}
\usepackage{fancyvrb}
\newcommand{\VerbBar}{|}
\newcommand{\VERB}{\Verb[commandchars=\\\{\}]}
\DefineVerbatimEnvironment{Highlighting}{Verbatim}{commandchars=\\\{\}}
% Add ',fontsize=\small' for more characters per line
\newenvironment{Shaded}{}{}
\newcommand{\AlertTok}[1]{\textcolor[rgb]{1.00,0.00,0.00}{\textbf{#1}}}
\newcommand{\AnnotationTok}[1]{\textcolor[rgb]{0.38,0.63,0.69}{\textbf{\textit{#1}}}}
\newcommand{\AttributeTok}[1]{\textcolor[rgb]{0.49,0.56,0.16}{#1}}
\newcommand{\BaseNTok}[1]{\textcolor[rgb]{0.25,0.63,0.44}{#1}}
\newcommand{\BuiltInTok}[1]{\textcolor[rgb]{0.00,0.50,0.00}{#1}}
\newcommand{\CharTok}[1]{\textcolor[rgb]{0.25,0.44,0.63}{#1}}
\newcommand{\CommentTok}[1]{\textcolor[rgb]{0.38,0.63,0.69}{\textit{#1}}}
\newcommand{\CommentVarTok}[1]{\textcolor[rgb]{0.38,0.63,0.69}{\textbf{\textit{#1}}}}
\newcommand{\ConstantTok}[1]{\textcolor[rgb]{0.53,0.00,0.00}{#1}}
\newcommand{\ControlFlowTok}[1]{\textcolor[rgb]{0.00,0.44,0.13}{\textbf{#1}}}
\newcommand{\DataTypeTok}[1]{\textcolor[rgb]{0.56,0.13,0.00}{#1}}
\newcommand{\DecValTok}[1]{\textcolor[rgb]{0.25,0.63,0.44}{#1}}
\newcommand{\DocumentationTok}[1]{\textcolor[rgb]{0.73,0.13,0.13}{\textit{#1}}}
\newcommand{\ErrorTok}[1]{\textcolor[rgb]{1.00,0.00,0.00}{\textbf{#1}}}
\newcommand{\ExtensionTok}[1]{#1}
\newcommand{\FloatTok}[1]{\textcolor[rgb]{0.25,0.63,0.44}{#1}}
\newcommand{\FunctionTok}[1]{\textcolor[rgb]{0.02,0.16,0.49}{#1}}
\newcommand{\ImportTok}[1]{\textcolor[rgb]{0.00,0.50,0.00}{\textbf{#1}}}
\newcommand{\InformationTok}[1]{\textcolor[rgb]{0.38,0.63,0.69}{\textbf{\textit{#1}}}}
\newcommand{\KeywordTok}[1]{\textcolor[rgb]{0.00,0.44,0.13}{\textbf{#1}}}
\newcommand{\NormalTok}[1]{#1}
\newcommand{\OperatorTok}[1]{\textcolor[rgb]{0.40,0.40,0.40}{#1}}
\newcommand{\OtherTok}[1]{\textcolor[rgb]{0.00,0.44,0.13}{#1}}
\newcommand{\PreprocessorTok}[1]{\textcolor[rgb]{0.74,0.48,0.00}{#1}}
\newcommand{\RegionMarkerTok}[1]{#1}
\newcommand{\SpecialCharTok}[1]{\textcolor[rgb]{0.25,0.44,0.63}{#1}}
\newcommand{\SpecialStringTok}[1]{\textcolor[rgb]{0.73,0.40,0.53}{#1}}
\newcommand{\StringTok}[1]{\textcolor[rgb]{0.25,0.44,0.63}{#1}}
\newcommand{\VariableTok}[1]{\textcolor[rgb]{0.10,0.09,0.49}{#1}}
\newcommand{\VerbatimStringTok}[1]{\textcolor[rgb]{0.25,0.44,0.63}{#1}}
\newcommand{\WarningTok}[1]{\textcolor[rgb]{0.38,0.63,0.69}{\textbf{\textit{#1}}}}
\setlength{\emergencystretch}{3em} % prevent overfull lines
\providecommand{\tightlist}{%
  \setlength{\itemsep}{0pt}\setlength{\parskip}{0pt}}
\setcounter{secnumdepth}{5}
\ifLuaTeX
  \usepackage{selnolig}  % disable illegal ligatures
\fi
\IfFileExists{bookmark.sty}{\usepackage{bookmark}}{\usepackage{hyperref}}
\IfFileExists{xurl.sty}{\usepackage{xurl}}{} % add URL line breaks if available
\urlstyle{same}
\hypersetup{
  colorlinks=true,
  linkcolor={blue},
  filecolor={Maroon},
  citecolor={Blue},
  urlcolor={blue},
  pdfcreator={LaTeX via pandoc}}

\author{}
\date{}

\begin{document}

\hypertarget{linear-algebra-ux69cbux9020ux304bux3089ux7406ux89e3ux3059ux308bux7ddaux5f62ux4ee3ux6570}{%
\section{Linear Algebra ---
構造から理解する線形代数}\label{linear-algebra-ux69cbux9020ux304bux3089ux7406ux89e3ux3059ux308bux7ddaux5f62ux4ee3ux6570}}

このリポジトリは、線形代数を「公式を覚える科目」ではなく、\textbf{空間・線形写像・分解・近似という一つの体系}として理解するための教材です。

\hypertarget{ux5b66ux7fd2ux306eux4e2dux5fc3ux3068ux306aux308bux554fux3044}{%
\subsection{学習の中心となる問い}\label{ux5b66ux7fd2ux306eux4e2dux5fc3ux3068ux306aux308bux554fux3044}}

線形代数では、次の問いを繰り返し扱います。

\begin{enumerate}
\def\labelenumi{\arabic{enumi}.}
\tightlist
\item
  どのようなベクトルを組み合わせると、どの空間を表現できるか。
\item
  線形写像は空間をどのように変形するか。
\item
  その変形によって、何が保存され、何が失われるか。
\item
  行列を適切な基底で見たとき、その本質はどこまで単純になるか。
\item
  厳密な解が存在しない場合、最も自然な近似解をどう定義するか。
\item
  固有値分解や SVD によって、線形変換の構造をどう読み解くか。
\end{enumerate}

\hypertarget{ux63a8ux5968ux5b66ux7fd2ux9806ux5e8f}{%
\subsection{推奨学習順序}\label{ux63a8ux5968ux5b66ux7fd2ux9806ux5e8f}}

\begin{Shaded}
\begin{Highlighting}[]
\NormalTok{ベクトル・ベクトル空間}
\NormalTok{        ↓}
\NormalTok{行列・線形写像}
\NormalTok{        ↓}
\NormalTok{連立一次方程式}
\NormalTok{        ↓}
\NormalTok{内積・直交・射影}
\NormalTok{        ↓}
\NormalTok{行列式}
\NormalTok{        ↓}
\NormalTok{固有値・固有ベクトル}
\NormalTok{        ↓}
\NormalTok{行列分解（LU / QR / EVD / SVD）}
\NormalTok{        ↓}
\NormalTok{最小二乗・擬似逆行列}
\NormalTok{        ↓}
\NormalTok{低ランク近似・PCA・画像圧縮・ノイズ除去}
\NormalTok{        ↓}
\NormalTok{数値線形代数・発展トピック}
\end{Highlighting}
\end{Shaded}

\hypertarget{ux30c7ux30a3ux30ecux30afux30c8ux30ea}{%
\subsection{ディレクトリ}\label{ux30c7ux30a3ux30ecux30afux30c8ux30ea}}

\begin{itemize}
\tightlist
\item
  \texttt{docs/}:
  解説。定義だけでなく、直感・導出・幾何学的意味・概念間の接続を重視します。
\item
  \texttt{exercises/}: 各章の演習問題。
\item
  \texttt{solutions/}: 演習の解答と考え方。
\item
  \texttt{examples/}: NumPy 実装、可視化、応用例。
\item
  \texttt{assets/}: 図や模式図を置くための領域。
\end{itemize}

まずは
\href{docs/00_overview/02_learning_path.pdf}{\texttt{docs/00\_overview/02\_learning\_path.md}}
から始めてください。

\end{document}
